\section{Introducción}

El transformador monofásico es una de las máquinas eléctricas más fundamentales dentro de los sistemas de potencia: permite elevar o reducir niveles de tensión mediante inducción electromagnética entre dos devanados acoplados por un núcleo ferromagnético común. Su versatilidad, eficiencia y relativa simplicidad constructiva lo convierten en un componente indispensable en las redes de transmisión y distribución, así como en innumerables aplicaciones industriales y domésticas~\cite{chapman,mit}.

En esta práctica de laboratorio se abordará la obtención del modelo circuital del transformador monofásico a partir de ensayos normalizados: medición de resistencia de devanados, determinación de la relación de transformación, verificación de polaridad, ensayo de cortocircuito y ensayo de vacío. Dichos ensayos permiten extraer los parámetros del circuito equivalente y constituyen la base para predecir el comportamiento del transformador bajo cualquier condición de carga.

\newpage

\section{Objetivos}

\subsection*{Objetivo General}
\begin{itemize}
    \item Determinar experimentalmente los parámetros del circuito equivalente de un transformador monofásico mediante ensayos normalizados, conforme a la norma COVENIN 3172-95.
\end{itemize}

\subsection*{Objetivos Específicos}
\begin{itemize}
    \item Medir la resistencia óhmica de los devanados primario y secundario por el método del voltímetro-amperímetro y referirla a la temperatura de operación.
    \item Determinar la relación de transformación del transformador mediante el método del voltímetro.
    \item Verificar la polaridad (aditiva o sustractiva) de los devanados empleando corriente alterna.
    \item Ejecutar el ensayo de cortocircuito para obtener las pérdidas en el cobre, la impedancia de cortocircuito y la tensión de cortocircuito.
    \item Ejecutar el ensayo de vacío para determinar las pérdidas en el hierro, la corriente de vacío, la resistencia de pérdidas en el núcleo y la reactancia de magnetización.
    \item Reforzar el uso de normas técnicas venezolanas e internacionales para la realización de pruebas en dispositivos eléctricos.
\end{itemize}

\newpage
